% Options for packages loaded elsewhere
\PassOptionsToPackage{unicode}{hyperref}
\PassOptionsToPackage{hyphens}{url}
\PassOptionsToPackage{dvipsnames,svgnames,x11names}{xcolor}
%
\documentclass[
]{article}
\usepackage{amsmath,amssymb}
\usepackage{iftex}
\ifPDFTeX
  \usepackage[T1]{fontenc}
  \usepackage[utf8]{inputenc}
  \usepackage{textcomp} % provide euro and other symbols
\else % if luatex or xetex
  \usepackage{unicode-math} % this also loads fontspec
  \defaultfontfeatures{Scale=MatchLowercase}
  \defaultfontfeatures[\rmfamily]{Ligatures=TeX,Scale=1}
\fi
\usepackage{lmodern}
\ifPDFTeX\else
  % xetex/luatex font selection
\fi
% Use upquote if available, for straight quotes in verbatim environments
\IfFileExists{upquote.sty}{\usepackage{upquote}}{}
\IfFileExists{microtype.sty}{% use microtype if available
  \usepackage[]{microtype}
  \UseMicrotypeSet[protrusion]{basicmath} % disable protrusion for tt fonts
}{}
\makeatletter
\@ifundefined{KOMAClassName}{% if non-KOMA class
  \IfFileExists{parskip.sty}{%
    \usepackage{parskip}
  }{% else
    \setlength{\parindent}{0pt}
    \setlength{\parskip}{6pt plus 2pt minus 1pt}}
}{% if KOMA class
  \KOMAoptions{parskip=half}}
\makeatother
\usepackage{xcolor}
\usepackage{color}
\usepackage{fancyvrb}
\newcommand{\VerbBar}{|}
\newcommand{\VERB}{\Verb[commandchars=\\\{\}]}
\DefineVerbatimEnvironment{Highlighting}{Verbatim}{commandchars=\\\{\}}
% Add ',fontsize=\small' for more characters per line
\usepackage{framed}
\definecolor{shadecolor}{RGB}{38,38,38}
\newenvironment{Shaded}{\begin{snugshade}}{\end{snugshade}}
\newcommand{\AlertTok}[1]{\textcolor[rgb]{1.00,0.37,0.37}{\textbf{#1}}}
\newcommand{\AnnotationTok}[1]{\textcolor[rgb]{0.35,0.35,0.35}{\textit{#1}}}
\newcommand{\AttributeTok}[1]{\textcolor[rgb]{0.00,0.53,0.69}{#1}}
\newcommand{\BaseNTok}[1]{\textcolor[rgb]{1.00,0.69,0.00}{#1}}
\newcommand{\BuiltInTok}[1]{\textcolor[rgb]{0.00,0.53,0.69}{#1}}
\newcommand{\CharTok}[1]{\textcolor[rgb]{1.00,0.84,0.00}{#1}}
\newcommand{\CommentTok}[1]{\textcolor[rgb]{0.35,0.35,0.35}{\textit{#1}}}
\newcommand{\CommentVarTok}[1]{\textcolor[rgb]{0.35,0.35,0.35}{\textit{#1}}}
\newcommand{\ConstantTok}[1]{\textcolor[rgb]{1.00,0.69,0.00}{#1}}
\newcommand{\ControlFlowTok}[1]{\textcolor[rgb]{0.69,0.53,1.00}{\textbf{#1}}}
\newcommand{\DataTypeTok}[1]{\textcolor[rgb]{0.00,0.53,0.69}{#1}}
\newcommand{\DecValTok}[1]{\textcolor[rgb]{1.00,0.69,0.00}{#1}}
\newcommand{\DocumentationTok}[1]{\textcolor[rgb]{0.35,0.35,0.35}{\textit{#1}}}
\newcommand{\ErrorTok}[1]{\textcolor[rgb]{1.00,0.37,0.37}{\textbf{#1}}}
\newcommand{\ExtensionTok}[1]{\textcolor[rgb]{0.00,0.53,0.69}{#1}}
\newcommand{\FloatTok}[1]{\textcolor[rgb]{1.00,0.69,0.00}{#1}}
\newcommand{\FunctionTok}[1]{\textcolor[rgb]{0.37,0.69,0.53}{#1}}
\newcommand{\ImportTok}[1]{\textcolor[rgb]{0.00,0.53,0.69}{#1}}
\newcommand{\InformationTok}[1]{\textcolor[rgb]{1.00,0.84,0.00}{#1}}
\newcommand{\KeywordTok}[1]{\textcolor[rgb]{0.69,0.53,1.00}{\textbf{#1}}}
\newcommand{\NormalTok}[1]{\textcolor[rgb]{0.74,0.74,0.74}{#1}}
\newcommand{\OperatorTok}[1]{\textcolor[rgb]{0.74,0.74,0.74}{#1}}
\newcommand{\OtherTok}[1]{\textcolor[rgb]{0.74,0.74,0.74}{#1}}
\newcommand{\PreprocessorTok}[1]{\textcolor[rgb]{1.00,0.84,0.00}{#1}}
\newcommand{\RegionMarkerTok}[1]{\textcolor[rgb]{0.74,0.74,0.74}{#1}}
\newcommand{\SpecialCharTok}[1]{\textcolor[rgb]{1.00,0.84,0.00}{#1}}
\newcommand{\SpecialStringTok}[1]{\textcolor[rgb]{1.00,0.84,0.00}{#1}}
\newcommand{\StringTok}[1]{\textcolor[rgb]{1.00,0.84,0.00}{#1}}
\newcommand{\VariableTok}[1]{\textcolor[rgb]{0.00,0.53,0.69}{#1}}
\newcommand{\VerbatimStringTok}[1]{\textcolor[rgb]{1.00,0.84,0.00}{#1}}
\newcommand{\WarningTok}[1]{\textcolor[rgb]{1.00,0.37,0.37}{\textbf{#1}}}
\usepackage{longtable,booktabs,array}
\usepackage{calc} % for calculating minipage widths
% Correct order of tables after \paragraph or \subparagraph
\usepackage{etoolbox}
\makeatletter
\patchcmd\longtable{\par}{\if@noskipsec\mbox{}\fi\par}{}{}
\makeatother
% Allow footnotes in longtable head/foot
\IfFileExists{footnotehyper.sty}{\usepackage{footnotehyper}}{\usepackage{footnote}}
\makesavenoteenv{longtable}
\usepackage{graphicx}
\makeatletter
\def\maxwidth{\ifdim\Gin@nat@width>\linewidth\linewidth\else\Gin@nat@width\fi}
\def\maxheight{\ifdim\Gin@nat@height>\textheight\textheight\else\Gin@nat@height\fi}
\makeatother
% Scale images if necessary, so that they will not overflow the page
% margins by default, and it is still possible to overwrite the defaults
% using explicit options in \includegraphics[width, height, ...]{}
\setkeys{Gin}{width=\maxwidth,height=\maxheight,keepaspectratio}
% Set default figure placement to htbp
\makeatletter
\def\fps@figure{htbp}
\makeatother
\usepackage{svg}
\ifLuaTeX
  \usepackage{luacolor}
  \usepackage[soul]{lua-ul}
\else
  \usepackage{soul}
\fi
\setlength{\emergencystretch}{3em} % prevent overfull lines
\providecommand{\tightlist}{%
  \setlength{\itemsep}{0pt}\setlength{\parskip}{0pt}}
\setcounter{secnumdepth}{-\maxdimen} % remove section numbering
% ---------------------------------------------------------------------------
% header.tex — pandoc MD→PDF code block styling
% Pair with: --highlight-style .fcc/pdf/p10k.theme
%
% Approach: tcolorbox owns the block background entirely.
% fvextra patches Highlighting directly to fix the ␣ bug.
% ---------------------------------------------------------------------------

% ---------------------------------------------------------------------------
% Page margins
% ---------------------------------------------------------------------------
\usepackage[top=2cm, bottom=2cm, left=2cm, right=2cm]{geometry}

% ---------------------------------------------------------------------------
% Hyperlinks — controlled via -V colorlinks=true passed to pandoc.
% pandoc's template conditionally sets colorlinks/hidelinks based on that
% variable — no hyperref config needed here.
% ---------------------------------------------------------------------------

% ---------------------------------------------------------------------------
% Images — required for \includegraphics in title page template
% ---------------------------------------------------------------------------
\usepackage{graphicx}
\usepackage{xcolor}
\usepackage{fvextra}
\usepackage{etoolbox}
\usepackage[most]{tcolorbox}

% ---------------------------------------------------------------------------
% p10k colours (xterm-256 exact hex)
% ---------------------------------------------------------------------------
\definecolor{termbg} {HTML}{262626}
\definecolor{termfg} {HTML}{bcbcbc}

% ---------------------------------------------------------------------------
% Fix ␣ bug: patch the Highlighting environment that pandoc defines via
% \DefineVerbatimEnvironment{Highlighting}{Verbatim}{commandchars=...}
% fvextra's \fvset doesn't reach it because the options on DefineVerbatimEnvironment
% take precedence. We redefine it after \begin{document} to inject breaklines.
% ---------------------------------------------------------------------------
\AtBeginDocument{%
  % Redefine Highlighting to add fvextra breaklines + correct foreground colour.
  % Preserve the commandchars option pandoc requires for token coloring.
  \DefineVerbatimEnvironment{Highlighting}{Verbatim}{%
    commandchars=\\\{\},%
    breaklines=true,%
    breakanywhere=false,%
    breaksymbolleft={},%
    fontsize=\small,%
  }%
  % Suppress hyphenation at font level — prevents TeX inserting hyphens on line breaks.
  \BeforeBeginEnvironment{Verbatim}{\hyphenchar\font=-1\relax}%
  \AfterEndEnvironment{Verbatim}{\hyphenchar\font=`\-\relax}%
}

% ---------------------------------------------------------------------------
% Redefine Shaded using tcolorbox — completely independent of framed/snugshade.
% \AtBeginDocument so it fires after pandoc defines Shaded at line ~47.
% ---------------------------------------------------------------------------
\AtBeginDocument{%
  \ifdef{\Shaded}{}{\newenvironment{Shaded}{}{}}%
  \renewenvironment{Shaded}{%
    \begin{tcolorbox}[%
      enhanced,%
      arc=3pt,%
      boxrule=0.5pt,%
      colback=termbg,%
      colframe=termbg!60!black,%
      coltext=termfg,%
      left=6pt, right=6pt, top=4pt, bottom=4pt,%
      before skip=8pt, after skip=8pt,%
      fontupper=\color{termfg}\small,%
    ]%
  }{%
    \end{tcolorbox}%
  }%
}

% ---------------------------------------------------------------------------
% Inline code styling
% Use \scalerel to allow the text to shrink when inside narrow table columns.
% \texttt globally uses \footnotesize so long identifiers fit in pandoc's
% hardcoded longtable column proportions.
% ---------------------------------------------------------------------------
\usepackage{relsize}
\let\oldtexttt\texttt
\renewcommand{\texttt}[1]{%
  \colorbox{termbg}{{\ttfamily\relsize{-1}\color{termfg}#1}}%
}

% ---------------------------------------------------------------------------
% Suppress all hyphenation in prose.
% \pretolerance=10000 makes TeX accept slightly loose lines before trying
% hyphenation. Combined with \hyphenpenalty=10000 it effectively disables
% hyphenation without breaking the layout engine entirely.
% ---------------------------------------------------------------------------
\pretolerance=10000
\hyphenpenalty=10000
\exhyphenpenalty=10000

% ---------------------------------------------------------------------------
% Table row colours — alternating light grey / white body rows, medium grey
% header. Applies to all longtables pandoc generates.
%
% \rowcolors{<start row>}{<odd colour>}{<even colour>}
% Start row 1 so the header row is row 0 (styled separately via \rowcolor).
% \arrayrulecolor keeps the rule lines visible against the light background.
% ---------------------------------------------------------------------------
\usepackage{colortbl}
\definecolor{tableheader}{HTML}{cccccc}
\definecolor{tablerowodd}{HTML}{f5f5f5}
\definecolor{tableroweven}{HTML}{ffffff}

\AtBeginDocument{%
  \rowcolors{2}{tablerowodd}{tableroweven}%
}

% Style the header row — pandoc wraps it in \endfirsthead / \endhead blocks.
% Redefine \toprule and \midrule to inject \rowcolor on the header row.
% \arrayrulecolor darkens the rule lines so they remain visible on light bg.
\arrayrulecolor{gray!50}

% Colour the header row by patching longtable's first-head block.
% \pretocmd injects \rowcolor{tableheader} immediately after \toprule,
% which is where pandoc places the header row.
\usepackage{longtable}
\AtBeginDocument{%
  \pretocmd{\toprule}{\rowcolor{tableheader}}{}{}%
}
\ifdefined\setmonofont
  \setmonofont{DejaVu Sans Mono}
\fi
\ifLuaTeX
  \usepackage{selnolig}  % disable illegal ligatures
\fi
\usepackage{bookmark}
\IfFileExists{xurl.sty}{\usepackage{xurl}}{} % add URL line breaks if available
\urlstyle{same}
\hypersetup{
  colorlinks=true,
  linkcolor={blue},
  filecolor={Maroon},
  citecolor={blue},
  urlcolor={blue},
  pdfcreator={LaTeX via pandoc}}

\author{}
\date{}

\begin{document}

\renewcommand*\contentsname{Contents}
{
\hypersetup{linkcolor=}
\setcounter{tocdepth}{2}
\tableofcontents
}
\section{Feature reference}\label{feature-reference}

Everything holo-convert has to carry through, in one document: inline
marks, lists, tables, code, maths, images, links, breaks - and the
characters the pre-passes rewrite: arrows →, dashes --- and ticks ✓, and
a wikilink to {[}{[}Other note\textbar a label{]}{]} or {[}{[}Plain
target\#a-heading{]}{]}.

\subsection{Text}\label{text}

Plain, \emph{emphasis}, \textbf{strong}, \textbf{\emph{both}},
\st{strike}, \texttt{inline\ code}, \textsuperscript{super},
\textsubscript{sub}, a footnote\footnote{The footnote, with
  \texttt{code}.}, an \href{https://example.org}{external link}, and an
\hyperref[tables]{internal link}. A \hyperref[nowhere-at-all]{broken
anchor} and a near miss to \hyperref[tabels]{the tables}.

\begin{quote}
A block quote, with \textbf{strong} text inside.
\end{quote}

\begin{itemize}
\tightlist
\item[$\square$]
  an open task
\item[$\boxtimes$]
  a done task
\end{itemize}

\begin{enumerate}
\def\labelenumi{\arabic{enumi}.}
\tightlist
\item
  first
\item
  second

  \begin{itemize}
  \tightlist
  \item
    nested bullet
  \item
    another
  \end{itemize}
\end{enumerate}

\begin{description}
\tightlist
\item[Term]
Its definition.
\end{description}

\begin{center}\rule{0.5\linewidth}{0.5pt}\end{center}

\subsection{Tables}\label{tables}

\begin{longtable}[]{@{}
  >{\raggedright\arraybackslash}p{(\columnwidth - 4\tabcolsep) * \real{0.2857}}
  >{\centering\arraybackslash}p{(\columnwidth - 4\tabcolsep) * \real{0.3810}}
  >{\raggedleft\arraybackslash}p{(\columnwidth - 4\tabcolsep) * \real{0.3333}}@{}}
\toprule\noalign{}
\begin{minipage}[b]{\linewidth}\raggedright
Left
\end{minipage} & \begin{minipage}[b]{\linewidth}\centering
Centre
\end{minipage} & \begin{minipage}[b]{\linewidth}\raggedleft
Right
\end{minipage} \\
\midrule\noalign{}
\endhead
\bottomrule\noalign{}
\endlastfoot
\begin{minipage}[t]{\linewidth}\raggedright
a\\
b\strut
\end{minipage} & \texttt{x} & 1.5 \\
a very long cell that has to wrap across the column width when rendered
& mid & 2 \\
\end{longtable}

\begin{longtable}[]{@{}
  >{\raggedright\arraybackslash}p{(\columnwidth - 2\tabcolsep) * \real{0.5000}}
  >{\raggedright\arraybackslash}p{(\columnwidth - 2\tabcolsep) * \real{0.5000}}@{}}
\toprule\noalign{}
\begin{minipage}[b]{\linewidth}\raggedright
Identifier
\end{minipage} & \begin{minipage}[b]{\linewidth}\raggedright
Meaning
\end{minipage} \\
\midrule\noalign{}
\endhead
\bottomrule\noalign{}
\endlastfoot
\texttt{converter-services-2026-10-10} & a long identifier in the first
column \\
\texttt{x} & short \\
\end{longtable}

\begin{center}\rule{0.5\linewidth}{0.5pt}\end{center}

\subsection{Code}\label{code}

\begin{Shaded}
\begin{Highlighting}[]
\KeywordTok{def}\NormalTok{ f(x):}
    \CommentTok{\# arrows → and dashes — and ticks ✓ stay as they are in code}
    \ControlFlowTok{return}\NormalTok{ x }\OperatorTok{**} \DecValTok{2}  \CommentTok{\# “quotes”, ≈}
\end{Highlighting}
\end{Shaded}

\begin{Shaded}
\begin{Highlighting}[]
\NormalTok{graph LR}
\NormalTok{  A {-}{-}\textgreater{} B}
\end{Highlighting}
\end{Shaded}

\subsection{Maths}\label{maths}

Inline \(e^{i\pi} + 1 = 0\) and display:

\[\sum_{k=1}^{n} k = \frac{n(n+1)}{2}\]

\subsection{Images}\label{images}

\begin{figure}
\centering
\includesvg{diagram.svg}
\caption{A vector diagram}
\end{figure}

\begin{figure}
\centering
\includegraphics{pixel.png}
\caption{A small picture}
\end{figure}

Line one\\
line two, after a hard break.

\newpage

\subsection{After a page break}\label{after-a-page-break}

The end.

\end{document}
